\subsubsection{Case 0,1}

The key of this case is the following identity

\begin{dmath}\label{q2-measure-Pi-sing}
	\tikzfig{q-2-2-measure-Pi-sing}
\end{dmath}

\begin{proof}
	We check the cases where $P = X, Y, $ and $Z$.
\begin{align*}
	\tikzfig{q-2-2.3}
\end{align*}
\end{proof}

Applying \eqref{q2-measure-Pi-sing} multiple times we can get the following identity.

\begin{dmath}\label{q2-measure-Pi-mult}
	\tikzfig{q-2-2-measure-Pi-mult}
\end{dmath}


I claim that $C_{0,1} = (X H) \otimes P$ works, where $P$ is the Pauli corresponding to the initial Pauli measurement.
Check the following diagram

\begin{align*}
	\tikzfig{q-2-2.2}
\end{align*}

Where in the last step we used the fact that $P = \bigotimes P_i$ and $PP = I$ (as $P$ is Pauli).



